{\large\textbf{E1: Hidden wire}}

\textbf{Experimental setup and tasks}

A very long copper wire runs horizontally at unknown depth $h$ under a
horizontal square surface of side length $L = 100.0\,\mathrm{mm}$. The sides of
the square are oriented West-East (the $x$-axis) and South-North (the
$y$-axis), as shown in the figure. The origin of the coordinate system coincides
with the South-West corner of the square.

\begin{center}\includegraphics[width=0.30\linewidth]{figures/EuPhO_2021_Q4_fig1.png}\end{center}

The wire is connected to an adjustable DC source (not shown in the figure),
which can provide a current $I$ in the range from $-5\,\mathrm{A}$ to
$5\,\mathrm{A}$. The reversal of the sign of the current corresponds to a
reversal of the polarity of the source. A small compass can be placed on the
square surface, (including its circumference) to sense the magnetic field of
the wire through the deflection angle $\varphi$ between the magnetic needle and
the North ($y$) direction. Positive $\varphi$ values correspond to an Eastward
deflection, as shown in the figure, while negative $\varphi$ correspond to a
Westward deflection. You can assume that:
\begin{itemize}
\item The magnetic needle is a point-like magnetic dipole, which can rotate
freely around the vertical axis, i.e. the compass is sensitive to the
horizontal component of the magnetic field only.
\item The height of the needle above the surface is negligible compared to the
depth of the wire beneath the surface, i.e. the needle is situated in the
$xy$-plane.
\end{itemize}

Design your experiment and make the necessary simulations to perform the
following tasks:

a. Determine the orientation of the wire with respect to the coordinate system
by specifying its equation in the form: $y = ax + b$, and estimate the
uncertainties of the parameters $a$ and $b$. Draw the wire position on a graph
and indicate the direction corresponding to a positive current $I$.

b. Determine the depth $h$ of the wire below the surface and the horizontal
component $B_\mathrm{E}$ of the Earth's magnetic field. In this task you are not
required to calculate the experimental uncertainties explicitly, however, your
final results must be represented with an appropriate number of significant
digits.

The magnetic permeability of free space is
\[
\mu_0 = 4\pi \times 10^{-7}\,\mathrm{T\,m/A}.
\]

\textbf{Description of the simulation software}

The command line program simulates the measurement of the deflection angle
$\varphi$ after providing the current $I$ and placing the compass at the
coordinates $x$ and $y$ on the surface.

A typical output of a single simulation cycle of the program looks like:
\begin{verbatim}
Enter I (A) between -5.0 and 5.0: 3.4
Enter X (mm) between 0 and 100: 55
Enter Y (mm) between 0 and 100: 31
PHI = -33 degrees
-------------------------------
Enter I (A) between -5.0 and 5.0: _
\end{verbatim}

First, you enter the current $I$ in $A$ (the number between $-5.0$ and $5.0$),
then coordinates $x$ and $y$ in $mm$ (the numbers between $0$ and $100$). Each
input is confirmed with the \textbf{Enter} key. The program will output the
value of $\varphi$ (\texttt{PHI}) in degrees (rounded to $1^\circ$) and return
to the initial prompt.

The current input $I$ will be rounded to $0.1\,\mathrm{A}$, the coordinate
inputs $x$, $y$ will be rounded to $1\,\mathrm{mm}$ before being used in
simulation. (There is no point in trying to input more precise numbers).

Every time you change the position of a compass, its real position used in
simulation differs from the input coordinates with an error about
$0.5\,\mathrm{mm}$. (It is a simulation of a limited precision when you place
an object).

Any time you need to quit the program, press \textbf{Ctrl+C}.
